%%
%% Copyright 2022 OXFORD UNIVERSITY PRESS
%%
%% This file is part of the 'oup-authoring-template Bundle'.
%% ---------------------------------------------
%%
%% It may be distributed under the conditions of the LaTeX Project Public
%% License, either version 1.2 of this license or (at your option) any
%% later version.  The latest version of this license is in
%%    http://www.latex-project.org/lppl.txt
%% and version 1.2 or later is part of all distributions of LaTeX
%% version 1999/12/01 or later.
%%
%% The list of all files belonging to the 'oup-authoring-template Bundle' is
%% given in the file `manifest.txt'.
%%
%% Template article for OXFORD UNIVERSITY PRESS's document class `oup-authoring-template'
%% with bibliographic references
%%

%%%CONTEMPORARY%%%
% \documentclass[unnumsec,webpdf,contemporary,large]{oup-authoring-template}%
%\documentclass[unnumsec,webpdf,contemporary,large,namedate]{oup-authoring-template}% uncomment this line for author year citations and comment the above
%\documentclass[unnumsec,webpdf,contemporary,medium]{oup-authoring-template}
%\documentclass[unnumsec,webpdf,contemporary,small]{oup-authoring-template}

%%%MODERN%%%
\documentclass[unnumsec,webpdf,modern,large]{oup-authoring-template}
%\documentclass[unnumsec,webpdf,modern,large,namedate]{oup-authoring-template}% uncomment this line for author year citations and comment the above
%\documentclass[unnumsec,webpdf,modern,medium]{oup-authoring-template}
%\documentclass[unnumsec,webpdf,modern,small]{oup-authoring-template}

%%%TRADITIONAL%%%
% \documentclass[unnumsec,webpdf,traditional,large]{oup-authoring-template}
%\documentclass[unnumsec,webpdf,traditional,large,namedate]{oup-authoring-template}% uncomment this line for author year citations and comment the above
%\documentclass[unnumsec,namedate,webpdf,traditional,medium]{oup-authoring-template}
%\documentclass[namedate,webpdf,traditional,small]{oup-authoring-template}

% \onecolumn % for one column layouts

% \usepackage{showframe}

% \graphicspath{{figures/}}

% line numbers
% \usepackage[mathlines, switch]{lineno}
% \usepackage[right]{lineno}

\theoremstyle{thmstyleone}%
\newtheorem{theorem}{Theorem}%  meant for continuous numbers
%%\newtheorem{theorem}{Theorem}[section]% meant for sectionwise numbers
%% optional argument [theorem] produces theorem numbering sequence instead of independent numbers for Proposition
\newtheorem{proposition}[theorem]{Proposition}%
%%\newtheorem{proposition}{Proposition}% to get separate numbers for theorem and proposition etc.
\theoremstyle{thmstyletwo}%
\newtheorem{example}{Example}%
\newtheorem{remark}{Remark}%
\theoremstyle{thmstylethree}%
\newtheorem{definition}{Definition}

%\newcommand{\overlaps}[1]{\textbf{overlaps}(\textbf{$#1$})}
\newcommand{\overl}[2]{\mathbf{ov}(#1, #2)}

\newcommand{\counts}[1]{\textbf{count}(\textbf{$#1$})}
\newcommand{\overbar}[1]{\mkern 1.5mu\overline{\mkern-1.5mu#1\mkern-1.5mu}\mkern 1.5mu}

\usepackage[english=british]{csquotes}
\usepackage{mathtools}
\usepackage{ulem}
\DeclarePairedDelimiter\size{\lvert}{\rvert}
\newcommand{\avglen}[1]{\mathbf{\overbar{\text{$#1$}}}}
\newcommand{\len}[1]{\textbf{length}(\text{$#1$})}
\newcommand{\gs}{\textbf{GS}}
\newcommand{\ot}{\textbf{OT}}
\newcommand{\lrge}{LRGE}
\newcommand{\relerr}{$\epsilon_{\text{rel}}$}
\newcommand{\rtwo}{$R^2$}
\newcommand{\ava}{\textit{ava}}
\newcommand{\twoset}{\textit{2set}}
\newcommand{\kmer}{\textit{k}-mer}
\newcommand{\changed}[1]{\textcolor{red}{#1}}
\newcommand{\newversion}[1]{\textcolor{blue}{#1}}


% ============================================================
% ============================================================
% ============================================================
% ============================================================

\begin{document}

\journaltitle{Bioinformatics}
\DOI{DOI HERE}
\copyrightyear{2024}
\pubyear{2024}
\access{Advance Access Publication Date: Day Month Year}
\appnotes{Paper}

\firstpage{1}
%\subtitle{Subject Section}

\title[Long read genome size estimation]{Genome size estimation from long read overlaps}

\author[1, $\ast$]{Michael B Hall\ORCID{0000-0003-3683-6208}}
\author[2]{Chenxi Zhou\ORCID{0000-0002-1735-2630}}
\author[1,3]{Lachlan J M Coin\ORCID{0000-0002-4300-455X}}

\authormark{Hall \& Coin}
\address[1]{\orgdiv{Department of Microbiology and Immunology}, \orgname{The University of Melbourne, at the Peter Doherty Institute for Infection and Immunity}, \orgaddress{\street{Melbourne}, \country{Australia}}}
\address[2]{\orgdiv{Department of Genetics}, \orgname{The University of Cambridge}, \orgaddress{\street{Cambridge}, \country{United Kingdom}}}
\address[3]{\orgdiv{Department of Clinical Pathology}, \orgname{The University of Melbourne}, \orgaddress{\street{Melbourne}, \country{Australia}}}

%Department of Genetics, University of Cambridge, Cambridge CB2 3EH, UK

\corresp[$\ast$]{Corresponding author. michael.hall2[at]unimelb.edu.au}

\received{Date}{0}{Year}
\revised{Date}{0}{Year}
\accepted{Date}{0}{Year}

%\editor{Associate Editor: Name}

\abstract{
\textbf{Summary:} Accurate genome size estimation is an important component of genomic analyses, though existing tools are primarily optimised for short-read data. We present \lrge{}, a novel tool that uses read-to-read overlap information to estimate genome size in a reference-free manner. \lrge{} calculates per-read genome size estimates by analysing the expected number of overlaps for each read, considering read lengths and a minimum overlap threshold. The final size is taken as the median of these estimates, ensuring robustness to outliers such as reads with no overlaps. Additionally, \lrge{} provides an expected confidence range for the estimate. 
 We validate \lrge{} on a large, diverse bacterial dataset and confirm it generalises to eukaryotic datasets. On bacterial genomes, \lrge{} outperforms \kmer{}-based methods in both accuracy and computational efficiency and produces genome size estimates comparable to those from assembly-based approaches, like Raven, while using significantly less computational resources.\\ 
\textbf{Availability and implementation:} Our method, \lrge{} (\textbf{L}ong \textbf{R}ead-based \textbf{G}enome size \textbf{E}stimation from overlaps), is implemented in Rust and is available as a precompiled binary for most architectures, a Bioconda package, a prebuilt container image, and a \href{https://crates.io/}{crates.io} package as a binary (\texttt{lrge}) or library (\texttt{liblrge}). The source code is available at \url{https://github.com/mbhall88/lrge} and an archive at \url{https://doi.org/10.5281/zenodo.17183812} under an MIT license. 
}
\keywords{bioinformatics, genome size, long read sequencing, read overlaps}


\maketitle

% ============================================================
% ============================================================

\section{Introduction}

Genome size is a fundamental genomic characteristic essential for downstream analyses such as genome assembly \citep{wick_benchmarking_2020} and evolutionary studies \citep{pellicer_plant_2020}. Accurate genome size estimation remains challenging, especially for non-model organisms and datasets with high heterogeneity or repetitive content. Existing methods primarily focus on short-read data and often require high computational resources or rely on pre-assembled references \citep{vurture_genomescope_2017,ranallo-benavidez_genomescope_2020}, limiting their applicability to modern long-read sequencing platforms from Pacific Biosciences (PacBio) and Oxford Nanopore Technologies (ONT).

The advancement of long-read sequencing technologies has made it relatively easy to generate high-quality bacterial genome assemblies \citep{wick_assembling_2023}. Combined with the ever-increasing throughput of sequencing, automated pipelines for tasks like variant calling and genome assembly are now common \citep{murigneux_micropipe_2021,petit_bactopia_2020,bouras_hybracter_2024}. These pipelines often require a genome size estimation from the user \citep{murigneux_micropipe_2021}, or optionally estimate one \citep{petit_bactopia_2020,bouras_hybracter_2024}. However, existing size estimation tools are typically designed for short reads \citep{ranallo-benavidez_genomescope_2020,vurture_genomescope_2017} and struggle with the higher error rates of long reads, which can introduce many erroneous \kmer{}s, the basis for some estimation methods \citep{ondov_mash_2016,ranallo-benavidez_genomescope_2020}. Additionally, determining sequencing depth requires knowledge of the genome size \citep{hall_rasusa_2022}, so pipelines that downsample sequencing data to improve computational efficiency must either rely on user-provided size estimates or calculate them automatically.

Here, we present a novel approach leveraging long-read overlap data obtained from minimap2 \citep{li_minimap2_2018} to provide accurate genome size estimates in a reference-free manner, without a reliance on \kmer{}s. By focusing on read-to-read overlaps, our method efficiently captures genome-wide patterns of coverage and redundancy, offering a robust alternative to \kmer{}-based and assembly-dependent techniques. This tool is designed for analysing long-read datasets, making genome size estimation accessible for a broader range of organisms and experimental contexts.

Our method, \lrge{}, estimates genome size by analysing how individual reads overlap with one another. For each query read, it calculates the expected number of overlaps with a set of target reads, considering both the lengths of the reads and a minimum overlap threshold. Using this overlap information, the method derives a genome size estimate for each query read. These per-read estimates are then aggregated using the median, ensuring robustness to outliers such as reads with no overlaps, which would otherwise produce infinite estimates. This approach leverages the statistical properties of read overlaps to infer genome size without requiring a pre-assembled reference.

% ============================================================
% ============================================================

\section{Methods}

Suppose the genome size is $\gs$. We assume that we have sequenced a set of $\size{T}$ target reads $T$ with average length $\overline{\ell}_T = \frac{1}{|T|} \sum_{t \in T} \ell(t)$, where $\ell(t)$ is the length of a specific read $t \in T$; as well as a second set of $\size{Q}$ query reads $Q$ with average read length $\overline{\ell}_Q$, defined in the same manner as $\overline{\ell}_T$.  We denote by $\overl{Q}{T}$ the set of reads in $Q$ which overlap with $T$ and by $T \setminus \{q\}$ the set of reads in $T$ not including $q$, which enables us to consider the case that $Q$ and $T$ may include the same reads.  We make the assumption that reads in $\overl{Q}{T}$  overlap due to originating from neighbouring regions in the genome and not due to repetitive DNA elements. To minimise the impact of repetitive DNA we provide an option to filter alignments due to internal matches \citep{li2016minimap}, for use with eukaryotic genomes. We consider internal matches to be those overlaps where one read is fully contained within another.

If the target read $t$ is located at a known position on the genome, how often would a randomly placed query read $q$, drawn from the $\gs - \ell(q)$ possible positions, be expected to overlap with $t$?
If we consider that the target read starts at genome position $s(t)$ and ends at $e(t)$ then to achieve an overlap of at least $\ot$ bases, the query read start position $s(q)$ must lie within the interval $$[s(t) +\ot - \ell(q), e(t) +1 -\ot].$$ The size of this interval is $\ell(q) +\ell(t) -2*\ot +1$.  We make an implicit assumption here that $s(t) \geq ( \ell(q)-\ot)$ and that $e(t) < \gs - (\ell(q)-\ot)$, which holds for the vast majority of reads $t$.  In the edge cases where this assumption does not hold, the size of the interval is reduced, with lower bound of $\ell(q) -\ot$, however we will ignore this edge effect from here onwards. The total number of available positions for the query read to start is $\gs - \ell(q)$. Thus the probability that a read $t$ overlaps with a distinct read $q$ is equal to
\begin{equation}
    P(\size{\overl{\{t\}}{\{q\}}} >0 ) =  \frac{\ell(q) + \ell(t) - 2 \cdot \ot +1 }{\gs-\ell(q)}
\end{equation}
where we approximate $\ot$ as the base-pair threshold on the minimal overlap length, as defined by the minimap2 minimum chaining score, \verb|-m|, (default 100bp for overlaps) \citep{li_minimap2_2018}. The expected count of overlaps is thus given by 
\begin{equation}
 \textbf{E}(\size{\overl{T \setminus \{q\} }{\{q\}}}) = \size{T \setminus \{q\}} \frac{\ell(q) + \overline{\ell}_{T \setminus \{q\}} - 2 \cdot \ot +1 }{\gs-\ell(q)}
\end{equation}
which gives us a per-read estimate of $\gs$, 
\begin{equation}\label{eq:3}
    \gs_{T,q} \approx \ell(q) + \size{T \setminus \{q\}} \frac{\ell(q) + \overline{\ell}_{T \setminus \{q\}} - 2 \cdot \ot +1 }{\size{\overl{T \setminus \{q\}}{\{q\}}}}
\end{equation}


We calculate the overall genome size estimate, $\gs_{T,Q}$, as the $\textbf{\text{median}}$ of the defined per-read estimates:
%\begin{equation}
%    \gs_{T,Q} = \textbf{\text{median}}\left( \{ \gs_{T, q} \mid q \in Q \& \size{\overl{T \setminus q}{\{q\}}} >0 \} \right)
%\end{equation}

%\changed{
%\begin{equation}
%    \gs_{T,Q} = \textbf{\text{median}}\left({q \in Q \wedge \size{\overl{T \setminus \%{q\}}{\{q\}}} >0}\left( \{ \gs_{T, q}  \}\right)\right)
%\end{equation}
%}


%% reviewer suggestion \gs_{T,Q} = median(\{GS_{T,q} | q \in Q \wedge |ov(T \setminus \{q\}, \{q\})| > 0\})
\newversion{
\begin{equation}
    \gs_{T,Q} = \textbf{\text{median}}\left( \{ \gs_{T, q} \mid q \in Q  \wedge \size{\overl{T \setminus \{q\}}{\{q\}}} >0   \} \right)
\end{equation}
}


as a read $q$ with no overlaps will return an undefined per-read estimate.
%Note that we can also calculate the total number of expected %overlaps between $Q \setminus T$ and $T$  as
%\begin{equation}
%\textbf{E}(\size{\overl{Q \setminus T}{T}})) = \size{Q \setminus T} %* \size{T} * \frac{\avglen{Q \setminus T} + \avglen{T} - 2*\ot}{\gs},
 %\end{equation}
 %which allows us to solve for $\gs$ as follows
%\begin{equation}
%    \gs_{\textbf{all}}(Q,T) \approx \size{Q \setminus T} * %\size{T}* \frac{\avglen{Q \setminus T} + \avglen{T} - 2*\ot}%{\size{\overl{Q \setminus T}{T}}}.
%\end{equation}

%Thus we have two potential estimates: $\gs_{\textbf{median}}$ and %$\gs_{\textbf{all}}$.

% E( \textbf{ overlaps S r_1} ) =

% ============================================================

\subsection{Implementation}

We have implemented this estimation process in the software tool \lrge{} using the Rust programming language, with overlaps being generated by minimap2 \citep{li_minimap2_2018}. \lrge{} offers a command-line application, along with a library API for use in other Rust projects.

The implementation provides two strategies for estimation: Two-set (\twoset{}), whereby $Q$ and $T$ are disjoint sets of reads; and all-vs-all (\ava{}), where $Q$ and $T$ are identical sets of reads.

The \twoset{} strategy has the advantage where $Q$ can be made smaller than $T$ to reduce the number of per-read estimates, $\gs_{T,q}$, as we ultimately take the median of these estimates. Additionally, it enables internally swapping $Q$ and $T$ so that the smaller dataset can serve as the reference for minimap2 alignment. While this does not affect the results, it significantly reduces the memory footprint during mapping when $T$ is much larger than $Q$. This strategy is the default used in \lrge{}, with $\size{Q} = 5000$ and $\size{T} = 10000$.

%In the \ava{} strategy, we ignore self-overlaps. 
%Both strategies count only the first overlap between each pair of reads, ignoring subsequent overlaps.

% ============================================================

\subsection{Evaluation}
\label{sec:eval}

We compared the genome size estimations from \lrge{} using the \ava{} and \twoset{} strategies, 
as well as GenomeScope2 \citep{ranallo-benavidez_genomescope_2020}, Mash \citep{ondov_mash_2016}, and Raven \citep{vaser_time-_2021}.

For \lrge{} \ava{} we used 25,000 randomly selected reads, while for the \twoset{} strategy, 
we used a target set size of 10,000 and a query set size of 5,000.

Mash (v2.3; \citep{ondov_mash_2016}) estimates were gathered using the subcommand \texttt{sketch} with the minimum \kmer{} copy number for filtering set to 10 and a sketch size of 100,000.

The parameters used for \lrge{} and Mash were selected based on a sweep of parameter combinations on a validation set described in the appendix (\nameref{app:param_explore}).

To generate GenomeScope2 estimates, we computed the histogram of \kmer{} frequencies using 
KMC (v3.2; \citep{kokot_kmc_2017}) with a \kmer{} size of 21, a minimum \kmer{} count (\texttt{-ci}) of 2, and a maximum \kmer{} count (\texttt{-cs}) of 10,000. These parameters were selected as they are the default recommended in the GenomeScope2 repository. We then ran GenomeScope2 (v2.0.1; \citep{ranallo-benavidez_genomescope_2020}) against the frequency histogram with the same settings, with ploidy of 1.

Raven is a genome assembly tool that is designed to be time- and memory-efficient and uses an overlap-layout-consensus approach \citep{vaser_time-_2021} - keeping with the overlap theme. We set Raven (v1.8.3) to perform no polishing (\verb|-p 0|) and use the size of the computed assembly as the \enquote{estimate}.

We evaluated the accuracy of each method using relative error, \relerr{}, which measures the percentage difference between the estimated value ($\hat{G}$) and the true value ($G$), scaled relative to the true value. It quantifies how close the estimate is to the true value, with positive values indicating overestimation and negative values indicating underestimation. It is defined as

\begin{equation}\label{eq:rel_err}
    \epsilon_{\text{rel}} = \left( \frac{\hat{G}}{G} - 1 \right) \cdot 100
\end{equation}

Statistical testing of the absolute relative error, CPU time, and maximum memory usage differences between methods were performed using a pairwise Tukey's range test. For the absolute relative error testing, we only compared between the same sequencing technology -- i.e., we do not compare PacBio for method A with ONT for method B.


\subsubsection{Prokaryote evaluation}
\label{sec:prok}

We performed a large-scale evaluation of the estimation methods using a set of 3370 publicly-available bacterial long-read sequencing runs -- 2468 ONT and 902 PacBio. These runs were selected as they were associated with a high-quality RefSeq assembly, which we took to be the true size. A detailed description of the data collection and filtering performed is provided in the appendix (\nameref{app:datasets}).
All experiments for prokaryotic samples were conducted on a server running Red Hat Enterprise Linux (RHEL) 9.4 equipped with Intel\textsuperscript{\tiny{\textregistered}} Xeon\textsuperscript{\tiny{\textregistered}} Gold 6448H CPUs. 

\subsubsection{Eukaryote evaluation}
\label{sec:methods_euk}

While the primary development focus for \lrge{} was bacterial genomes, we investigated how the method generalises to eukaryotes, such as multichromosome and diploid organisms. We used ONT reads from four model organisms with a high-quality RefSeq genome: \textit{Saccharomyces cerevisiae}, \textit{Drosophila melanogaster}, \textit{Arabidopsis thaliana} and \textit{Homo sapiens}.

Due to the larger genome sizes of these organisms, we used 10,000 query reads and 20,000 target reads for the \twoset{} strategy on \textit{S. cerevisiae} and default (25,000) for the \ava{} strategy. For \textit{D. melanogaster} and \textit{A. thaliana}, we used 100,000 target reads and 50,000 query reads for the \lrge{} \twoset{} approach and 100,000 reads for the \ava{} strategy. For \textit{H. sapiens}, we used 100,000 query and 2,000,000 target reads for the \twoset{} strategy and did not run \ava{} strategy due to memory constraints. We enabled the \texttt{-F} option to filter out internal matches (overlaps where one read is fully contained within another) and \texttt{--use-min-ref} option to allow the smaller $Q$ dataset to be used as reference for minimap2 alignment across all these eukaryotic datasets. The parameters were kept the same as for the prokaryote evaluation for all other methods. All experiments for eukaryote samples were conducted on a server running Ubuntu 22.04, equipped with a 2.90GHz Intel\textsuperscript{\tiny{\textregistered}} Xeon\textsuperscript{\tiny{\textregistered}} Gold 6226R CPU, with the system recognizing up to 64 logical processors.

% ============================================================
% ============================================================

\section{Results}

To evaluate the accuracy of the estimations produced by \lrge{}, we collected publicly-available bacterial long-read sequencing runs with associated high-quality reference assemblies (see \nameref{sec:prok}). We benchmarked the two \lrge{} strategies against three other methods: Mash, which uses sketching of \kmer{}s to produce an estimate; GenomeScope2, which performs statistical analyses on \kmer{} frequency spectra; and Raven, which is a rapid genome assembly method (see \nameref{sec:eval}).

Our primary metric of comparison is relative error (\relerr{}; \autoref{eq:rel_err}), which measures the percentage difference between the estimate and the true value, scaled by the true value. For example, a \relerr{} of 50\% for a sample with a genome size of 2 Mbp indicates the estimate is over by half the true size -- 3 Mbp. To simplify the evaluation, we use the absolute value of \relerr{} to make for easier comparison across samples.

\autoref{fig:abs_rel_err} shows the results of this comparison. The first observation is that the two \lrge{} strategies provide very similar estimates, with no significant difference. Second, \lrge{} and Raven perform better with ONT data than PacBio, while Mash and GenomeScope2 are the opposite. Raven provides the best genome size estimates, being statistically significantly better than all methods on PacBio data (mean: 3.5\%, median: 2.6\%), and GenomeScope2 and Mash on ONT data (mean: 2.9\%, median: 1.1\%). \lrge{} had significantly lower \relerr{} compared to GenomeScope2 and Mash for ONT data (\ava{}: mean 9.2\%, median 4.6\%; \twoset{}: mean 8.9\%, median 4.8\%), though the inverse was true for PacBio data. The full results can be seen in Supplementary Table S1.

A key observation is that \lrge{} performs substantially better on ONT data than on PacBio. This is likely due to lower read quality in the available PacBio datasets compared to ONT, as shown in \autoref{fig:read_len_qual}. Notably, many PacBio reads have quality scores of zero—likely placeholders resulting from FASTA-to-FASTQ conversion during submission to SRA/ENA. Additionally, both read length and quality significantly affect estimation accuracy, with longer and higher-quality reads leading to lower \relerr{} (\autoref{fig:relerr_len_qual}). For reads of equivalent length and quality, \relerr{} tends to be lower for ONT data than for PacBio. This suggests that read quality alone does not fully account for the performance gap, and we were unable to determine the underlying cause of this difference.

% ============================================================

\subsection{Outliers}
\label{sec:outliers}

GenomeScope2 and Raven tended to underestimate the genome size (\autoref{fig:rel_err}) on ONT data, though the median was approximately -1\% for Raven. In terms of outliers for \lrge{} there were two common patterns. Major underestimates were found to be cases where there was dramatically disproportionate depth across the genomes. In one example, a \textit{Pandoraea fibrosis} run (SRR9733840), the ribosomal RNA genes had up to 100,000x read depth. In another, a \textit{Enterobacter ludwigii} run (SRR12247681) had two plasmids with 50,000x and 160,000x depth. This dramatic difference in copy number for plasmids is a known mechanism of transient antibiotic heteroresistance \citep{nicoloff_three_2024}. These dramatic depth differences lead to underestimates because, when randomly selecting reads, there is a high likelihood that mostly plasmid or gene reads will be selected, and thus the estimate will reflect the (much smaller) size of the plasmid(s) or gene(s).

The other major outlier pattern was low quality runs leading to large overestimates. This is illustrated in \autoref{fig:qual} where we see runs with a \relerr{} over 50\% with average read qualities much lower than other samples. Assumably this lower quality makes it very difficult to overlap reads, leading to fewer overlaps, and thus larger genome size estimates.

\begin{figure}[t]
\centering
\includegraphics[width=0.5\textwidth]{method_absolute_relative_error.pdf}
\caption{Absolute relative error (y-axis) for each method's (x-axis) genome size estimation on ONT (black) and PacBio (orange) data. The y-axis is scaled according to a symmetric logarithm, which is linear between -1 and 1 and logarithmic (base 10) thereafter. The statistical annotations are the result of a Tukey's range test and are coloured by the sequencing platform being compared. Pairs with no annotation indicate no significant difference. The dashed lines in the violins are the quartiles.\label{fig:abs_rel_err}}
\end{figure}

% ============================================================

\subsection{Calibration of estimate confidence range}

As we calculate a genome size estimate for each read, we can determine a range within which the estimate is likely to fall with a given level of confidence. To achieve this, we scanned all percentile ranges of width 50\%, ensuring that the overall estimate (median) was contained within the range. \autoref{fig:perc_rng} illustrates this process and shows that the percentile range that maximises the proportion of samples with the true genome size within the corresponding range is 15–65\% for both \lrge{} strategies. We found that 92.5\% (\twoset{}) and 87.8\% (\ava{}) of samples had their true genome size fall within this estimate range.

On the prokaryote dataset, this range gave an interquantile range (IQR) of 0.53 and 0.44 relative size ($\frac{\text{estimate}}{\text{true size}}$) for the \twoset{} and \ava{} strategies, respectively. For example, an IQR of 0.53 means that for the \twoset{} strategy, the middle 50\% of estimates fall within a range that is 53\% of the true genome size, while for the \ava{} strategy, this range is only 44\%. This illustrates that the \ava{} strategy provides a more precise estimate, with less variability in the middle half of the predictions compared to the \twoset{} strategy.

% ============================================================

\subsection{Runtime and memory benchmark}

\autoref{fig:comp} shows the benchmark of CPU time and maximum memory usage for all methods. In terms of CPU time, \lrge{} \twoset{} (mean: 38.2s, median: 17.4s) and GenomeScope2 (mean: 31.9s, median: 31.6s) had the quickest runtime, with no significant difference between the two. Raven was significantly slower than all methods except \lrge{} \ava{}. The longest recorded CPU time, 2.4 hours, was \lrge{} \ava{}. Both \lrge{} strategies had large CPU time outliers, highlighted by the differences between their mean (\ava{}: 194.7s, \twoset{}: 38.2s) and median (\ava{}: 77.2s, \twoset{}: 17.4s). These outliers were enriched in samples with large underestimates (\autoref{fig:cpu_rel_err}). Together with the fact that samples with large overestimates tended to have very fast runtimes (\autoref{fig:cpu_rel_err}), this indicates that the runtime is proportional to the number of overlaps (see \nameref{sec:outliers}).

\lrge{} \twoset{} memory usage (mean: 776 MB, median: 749 MB) was significantly lower than all other methods, while Raven had the highest memory usage (mean: 7.2 GB, median: 7.7 GB). The highest recorded memory usage was Mash, with a maximum of 42 GB.  

The full results can be seen in Supplementary Table S2.

\begin{figure}[t]
\centering
\includegraphics[width=0.5\textwidth]{method_cpu_memory.pdf}
\caption{CPU time (left y-axis) and maximum memory usage (right y-axis) for each method (colours). The y-axis is log-scaled (base 10). The statistical annotations are the result of a Tukey's range test and are coloured by the method with the lower mean value. The dashed lines in the violins are the quartiles.\label{fig:comp}}
\end{figure}

% ============================================================

\subsection{Eukaryotes}

While the primary focus and development of this work has been bacteria-focused, we investigated how \lrge{} generalises to multi-chromosome and diploid organisms. We evaluated \lrge{}, along with the other three methods, against four model organisms: \textit{S. cerevisiae}, \textit{D. melanogaster}, \textit{A. thaliana}, , and \textit{Homo sapiens} (see \nameref{sec:methods_euk}). 
These results are summarised in \autoref{tab:euk} with full details provided in Supplementary Table S3. They demonstrate that \lrge{} generalises well to eukaryotic organisms, with all estimates closely aligning with the expected genome sizes and D. melanogaster recording the lowest \relerr{}. For \textit{A. thaliana}, we used the TAIR10 genome assembly as the reference; however, this version likely underrepresents the genome's repeat content and underestimates its true size. A recent pangenome study reports \textit{A. thaliana} genome sizes ranging from 128 to 148 Mb, with an average of approximately 135 Mb \citep{lian2024pan}. Based on this value, \lrge{} actually produces the most accurate estimate. Notably, \lrge{} consistently yields estimates within the range of the other three methods, highlighting its robustness.
%having the lowest \relerr{} for \textit{S. cerevisiae} and \textit{A. thaliana}, and a \relerr{} within 3-8\% on \textit{D. melanogaster}. 


\begin{table*}[t]
\caption{Estimates on eukaryote samples.\label{tab:euk}}
\begin{tabular*}{\textwidth}{@{}lccrlrrrr@{}}
\toprule
Organism                                  & Ploidy             & Chromosomes          & True size\textsuperscript{$\ast$}                   & Method      & Estimate\textsuperscript{$\ast$}    & \relerr{}\textsuperscript{$\dagger$} & Time\textsuperscript{$\ddagger$} & Memory (GB) \\ \midrule
\multirow{5}{*}{\textit{Saccharomyces cerevisiae}} & \multirow{5}{*}{1} & \multirow{5}{*}{16} & \multirow{5}{*}{12.2 Mbp}  & \lrge{} \twoset{}   & 13.8 Mbp  & 13.76          & 0:21:56  & 4.7              \\
                                                   &                    &                     &                              & \lrge{} \ava{}    & 14.1 Mbp  & 15.73          & 0:12:05  & 5.3              \\
                                                   &                    &                     &                              & Mash              & 55.7 Mbp  & 358.48         & \textbf{0:08:13}  & \textbf{0.9}             \\
                                                   &                    &                     &                              & GenomeScope2      & 15.1 Mbp  & 23.96          & 0:17:09  & 11.2             \\
                                                   &                    &                     &                              & Raven             & 13.0 Mbp  & \textbf{7.14}          & 3:20:32 & 51.3             \\ \midrule
\multirow{5}{*}{\textit{Drosophila melanogaster}}  & \multirow{5}{*}{2} & \multirow{5}{*}{7}  & \multirow{5}{*}{143.7 Mbp} & \lrge{} \twoset{}   & 144.8 Mbp & \textbf{0.77}          & 1:47:05  & 16.7              \\
                                                   &                    &                     &                              & \lrge{} \ava{}    & 145.9 Mbp & 1.55          & 1:25:55  & 21.3              \\
                                                   &                    &                     &                              & Mash              & 135.6 Mbp & -5.64          & \textbf{0:07:27}  & \textbf{1.8}              \\
                                                   &                    &                     &                              & GenomeScope2      & 150.6 Mbp & 4.76         & 0:15:27  & 11.2             \\
                                                   &                    &                     &                              & Raven             & 147.7 Mbp & 2.79          & 1:45:58 & 51.3             \\ \midrule
\multirow{5}{*}{\textit{Arabidopsis thaliana}}     & \multirow{5}{*}{2} & \multirow{5}{*}{5}  & \multirow{5}{*}{119.7 Mbp} & \lrge{} \twoset{}   & 143.1 Mbp & 19.55          & 12:00:15  & 29.4              \\
                                                   &                    &                     &                              & \lrge{} \ava{}    & 143.5 Mbp & 19.93          & 9:42:23  & 44.2              \\
                                                   &                    &                     &                              & Mash              & 147.2 Mbp & 22.96         & \textbf{0:16:02}  & \textbf{1.8}              \\
                                                   &                    &                     &                              & GenomeScope2      & 152.9 Mbp & 27.74          & 0:27:45  & 11.0             \\
                                                   &                    &                     &                              & Raven             & 123.5 Mbp & \textbf{3.16}          & 7:09:10 & 56.8             \\ \midrule
\multirow{5}{*}{\textit{Homo sapiens}}     & \multirow{5}{*}{2} & \multirow{5}{*}{24}  & \multirow{5}{*}{3117.3 Mbp} & \lrge{} \twoset{}   & 2807.8 Mbp & -9.93          & 194:34:48  & 46.3              \\
                                                   &                    &                     &                              & \lrge{} \ava{}    & - & -          & -  & -              \\
                                                   &                    &                     &                              & Mash              & 2307.0 Mbp & -25.99         & \textbf{1:24:03}  & \textbf{1.8}              \\
                                                   &                    &                     &                              & GenomeScope2      & 2962.3 Mbp & \textbf{-4.97}          & 2:55:47  & 11.3             \\
                                                   &                    &                     &                              & Raven             & 2900.5 Mbp & -6.95          & 365:38:27 & 112.6             \\ \botrule
\end{tabular*}
\begin{tablenotes}%
\item \textsuperscript{$\ast$} Size values have been rounded to megabase pair with one decimal place for compactness.
\item \textsuperscript{$\dagger$} \relerr{} is relative error expressed as a percentage, as defined in \autoref{eq:rel_err}.
\item \textsuperscript{$\ddagger$} CPU time formatted as [h]:mm:ss.
\item \textbf{Bold text} indicates the best value for the relevant metric for each organism.
\end{tablenotes}
\end{table*}


% ============================================================
% ============================================================

\section{Conclusion}

\lrge{} provides a robust and efficient method for genome size estimation tailored to long-read sequencing technologies. By leveraging read-to-read overlap data, it circumvents the limitations of \kmer{}-based methods and performs well across diverse datasets, including bacteria and eukaryotes. \lrge{} outperforms existing \kmer{}-based tools on bacterial genomes in both accuracy and computational efficiency and delivers estimates comparable in accuracy to those derived from an assembly-based approach, Raven, but with significantly reduced computational resource requirements. In practice, we find that excluding short or low-quality reads prior to running \lrge{} can improve estimation accuracy. Despite not being specifically designed for large, repetitive eukaryotic genomes, \lrge{} generalises well to these datasets, demonstrating both accuracy and robustness in its estimates. Its flexibility and scalability make \lrge{} a valuable addition to the bioinformatics toolkit, supporting applications in genome assembly, evolutionary studies, and sequencing depth estimation.

% ============================================================
% ============================================================

\section{Competing interests}

No competing interest is declared.

% ============================================================
% ============================================================

\section{Author contributions statement}

\textbf{Michael B Hall:} Conceptualization, Data curation, Formal Analysis, Investigation, Methodology, Software, Writing – original draft, Visualization, Writing – review \& editing.
\textbf{Chenxi Zhou:} Formal Analysis, Methodology, Software, Writing -review \& editing.
\textbf{Lachlan J M Coin:} Conceptualization, Formal Analysis, Funding acquisition, Investigation, Methodology, Software, Supervision, Writing – original draft, Writing – review \& editing. 

% ============================================================
% ============================================================

\section{Acknowledgments}

This research was supported by the University of Melbourne’s Research Computing Services and the Petascale Campus Initiative. We thank Ryan Wick for insightful suggestions and Leah Roberts for discussions relating to large plasmid copy numbers.

\subsection{Funding}

This work was supported, in part, by the Australian Government Medical Research Future Fund (MRFF) Genomics Health Futures Mission (GHFM) Flagships – Pathogen Genomics Grant (FSPGN000045) META-GP: DELIVERING A CLINICAL METAGENOMICS PLATFORM FOR AUSTRALIA. The funding body had no role in the design, analysis, interpretation, or writing of this work. 

% ============================================================
% ============================================================
\section{Data availability}

All code and metadata required to perform the analyses in this work are available on GitHub at \url{https://github.com/mbhall88/lrge} and are archived at Zenodo \citep{hall_2025_17183812} or in Supplementary Table S3 for the eukaryote samples.

% ============================================================
% ============================================================

\bibliographystyle{abbrvnat}
\bibliography{reference}

% ============================================================
% ============================================================
% ============================================================
% ============================================================
\begin{appendices}
% =====
% reset the figure counter and add an S prefix to the numbers
\renewcommand{\thefigure}{S\arabic{figure}}
\setcounter{figure}{0}
% =====

\section{Parameter exploration}\label{app:param_explore}

We explored the best combination of parameters to be used for the larger dataset in the main text, by validating on a much smaller dataset. For this, we used ONT reads from 14 different bacterial species \citep{hall_benchmarking_2024}. In particular, these are R10.4.1 simplex reads basecalled with the v4.3.0 Dorado super accurate (sup) model \citep{dorado}. 

% ============================================================

\subsection{\lrge{} two-set overlaps}\label{app:param_lrge_2set}

For the \lrge{} two-set (\twoset{}) approach, the three parameters we explore are the number of query and target reads used for the overlaps, $Q$ and $T$, respectively, and the read selection strategy. For the \textit{long} strategy, we select the $\size{Q}$ longest reads, $Q$ and the next $\size{T}$ longest reads, $T$. In the \textit{rand} strategy, we randomly select $Q$ and $T$ reads, ensuring these sets are disjoint, using Rasusa (v2.1.0) \citep{hall_rasusa_2022}. We then determine the overlaps using minimap2 (v2.28) \citep{li_minimap2_2018} with the $Q$ reads as the query and $T$ as the target. For the \textit{rand} strategy, we performed three replicates for each sample-$Q$-$T$ combination.

The metrics we used to assess the performance of the different values of $Q$ and $T$ are relative error (\autoref{eq:rel_err}; \relerr{}) and the coefficient of determination ($R^2$), with respect to the identity line, for each parameter combination. From \autoref{fig:2set-explore} we can see clearly that the \textit{long} strategy is inferior to the \textit{rand} strategy. $Q=5000$ with $T=10000$ had the best balance of $R^2$ and \relerr{} and were therefore selected as the default.

% ============================================================

\subsection{\lrge{} all-vs-all overlaps}\label{app:param_lrge_ava}

For the \lrge{} all-vs-all (\ava{}) approach, the two parameters we explore are the number of reads used for the overlaps, $n$, and the read selection strategy. For the \textit{long} strategy, we select the $n$ longest reads and in the \textit{rand} strategy, we select $n$ reads at random using Rasusa (v2.1.0) \citep{hall_rasusa_2022}. We then determine the overlaps using minimap2 (v2.28) \citep{li_minimap2_2018} with the $n$ reads as both the target and query. For the \textit{rand} strategy, we performed three replicates for each sample-$n$ combination.

From \autoref{fig:ava-explore} we can see clearly that the \textit{long} strategy is inferior to the \textit{rand} strategy. $n=10000$ had the highest $R^2$ and closest \relerr{} to zero, though $n=25000$ was very close. We chose to use 25000 as the default because of the fact it was so close to $n=10000$ and the fact that it performed much better on the \textit{long} strategy.


% ============================================================

\subsection{Mash}\label{app:param_mash}

The parameters we explored for Mash \citep{ondov_mash_2016} were the sketch size (\verb|-s|) and the minimum copies of a \kmer{} required to pass the noise filter (\verb|-m|). We ran the Mash subcommand \texttt{sketch} with these parameters for each of the 14 samples. \autoref{fig:mash-explore} shows the results of this parameter exploration, with a sketch size of 100000 and minimum copies of 10 giving the tightest distribution of relative error near zero.


% ============================================================
% ============================================================

\section{Dataset selection}\label{app:datasets}%

Metadata for all complete bacterial assemblies in NCBI's RefSeq (excluding MAGs) were downloaded using NCBI's \texttt{datasets} command line tool. We selected assemblies where: (i) the sequencing technology section mentions PacBio or Oxford Nanopore Technologies, (ii) the assembly release date was after 01/01/2016, (iii) CheckM contamination was below 3\%, (iv) the CheckM completeness was greater than 95\%, (v) the CheckM completeness percentile for the species was above 75\%. Next, we gathered metadata on the sequencing runs associated with the BioSample accession 
for each assembly. We removed assemblies which did not actually have any long read ONT or PacBio sequencing, despite criteria (i) from above. 

The long read data left after these filtering steps were downloaded using Kingfisher (v0.4; \citep{woodcroft_kingfisher_2024}). We then removed a further 1,130 sequencing runs due to low volume of sequencing reads, indicated by the readset not having sufficient reads to run \lrge{} \twoset{}, or GenomeScope2 giving an error stating it was unable to converge on an estimate (an error that typically indicates low sequencing depth). We also removed 22 runs which were actually Illumina reads, despite being labeled as PacBio or ONT on the NCBI SRA, and 16 runs for which a valid FTP address could not be found to download the data.

In the end, we were left with 3,370 long read sequencing runs.

% ============================================================
% ============================================================

\begin{figure*}[ht]
\centering
\includegraphics[width=\textwidth]{lrge_2set_param_explore.pdf}
\caption{\lrge{} two-set parameter exploration. Columns are the query read set size $Q$, while the rows are the target read set size $T$. Each subplot shows the estimated genome size (y-axis) against the true genome size (x-axis), with points coloured by the strategy. \relerr{} is the relative error and $R^2$ is the coefficient of determination. The subplot with the green border is the one that was selected as the default parameters.\label{fig:2set-explore}}
\end{figure*}

% ============================================================
% ============================================================

\begin{figure}[ht]
\centering
\includegraphics[width=0.5\textwidth]{lrge_ava_param_explore.pdf}
\caption{\lrge{} all-vs-all parameter exploration. Each subplot represents the number of reads ($N$) used and shows the estimated genome size (y-axis) against the true genome size (x-axis), with points coloured by the strategy. \relerr{} is the relative error and $R^2$ is the coefficient of determination. The subplot with the green border is the one that was selected as the default parameters.\label{fig:ava-explore}}
\end{figure}

% ============================================================
% ============================================================

\begin{figure}[ht]
\centering
\includegraphics[width=0.5\textwidth]{mash_param_explore.pdf}
\caption{Exploring the best-performing combination of Mash sketch size (x-axis) and minimum required copies of a \kmer{} to pass noise filtering (colours). The darker horizontal dashed line at $y=0$ indicates the optimal relative error. The y-axis is scaled according to a symmetric logarithm, which is linear between -1 and 1 and logarithmic (base 10) thereafter.\label{fig:mash-explore}}
\end{figure}

% ============================================================
% ============================================================

\begin{figure}[ht]
\centering
\includegraphics[width=0.5\textwidth]{platform_relative_error.pdf}
\caption{Relative error (y-axis) for each method's (x-axis) genome size estimation on ONT (black) and PacBio (orange) data. The y-axis is scaled according to a symmetric logarithm, which is linear between -1 and 1 and logarithmic (base 10) thereafter. The dashed lines in the violins are the quartiles.\label{fig:rel_err}}
\end{figure}

% ============================================================
% ============================================================

\begin{figure}[ht]
\centering
\includegraphics[width=0.5\textwidth]{avg_qual_relative_error.pdf}
\caption{Distribution of average read quality for samples with a relative error (\relerr{}) greater than 50\% (black) and all other samples (orange).\label{fig:qual}}
\end{figure}

% ============================================================
% ============================================================

\begin{figure}[ht]
\centering
\includegraphics[width=0.5\textwidth]{read_len_qual.pdf}
\caption{Distribution of read length and quality for all reads used by \lrge{} on ONT (blue) and PacBio (red) data.\label{fig:read_len_qual}}
\end{figure}


% ============================================================
% ============================================================

\begin{figure}[ht]
\centering
\includegraphics[width=0.5\textwidth]{relerr_vs_len_qual.pdf}
\caption{Impact of read length (left) and quality (right) on \lrge{} genome size estimation accuracy. Points and lines are coloured by sequencing platform. The lines indicate a linear regression fit to the points with the significance of the fit and the coefficient of determination ($R^2$) in the legend. Infinite estimates are those reads which have no overlaps.\label{fig:relerr_len_qual}}
\end{figure}

% ============================================================
% ============================================================

\begin{figure}[ht]
\centering
\includegraphics[width=0.5\textwidth]{percentile_range_scan.pdf}
\caption{Calibrating the optimal percentile range for \lrge{} estimates. The x-axis represents the lower end of the range. So an x value of 10 means a percentile range of 10-60. The y-axis represents the proportion of samples for which the true genome size lies within the corresponding percentile range on the x-axis. The two colours represent the two \lrge{} strategies.\label{fig:perc_rng}}
\end{figure}

% ============================================================
% ============================================================

\begin{figure}[ht]
\centering
\includegraphics[width=0.5\textwidth]{cpu_time_relative_error.pdf}
\caption{Relationship between CPU time (x-axis) and relative error (y-axis) for \lrge{} strategies (colours). The line represents a linear regression model fit to the data.\label{fig:cpu_rel_err}}
\end{figure}

\end{appendices}

\end{document}
